\documentclass[12pt]{article}

% Idioma y codificación
\usepackage[spanish, es-tabla]{babel}       %es-tabla para que se titule "Tabla"
\usepackage[utf8]{inputenc}

% Márgenes
\usepackage[a4paper,top=3cm,bottom=2.5cm,left=3cm,right=3cm]{geometry}

% Comentarios de bloque
\usepackage{verbatim}

% Paquetes de links
\usepackage[hidelinks]{hyperref}    % Permite enlaces
\usepackage{url}                    % redirecciona a la web

% Más opciones para enumeraciones
\usepackage{enumitem}


%------------------------------------------------------------------------

%Paquetes de figuras
\usepackage{caption}
\usepackage{subcaption} % Figuras al lado de otras
\usepackage{float}      % Poner figuras en el sitio indicado H.


% Paquetes de imágenes
\usepackage{graphicx}       % Paquete para añadir imágenes
\usepackage{transparent}    % Para manejar la opacidad de las figuras

% Paquete para usar colores
\usepackage[dvipsnames]{xcolor}
\usepackage{pagecolor}      % Para cambiar el color de la página

% Habilita tamaños de fuente mayores
\usepackage{fix-cm}


%------------------------------------------------------------------------

% Paquetes de matemáticas
\usepackage{amsmath, amsfonts, amssymb, mathrsfs}
\usepackage{cancel}     % Simplificar tachando
\usepackage{polynom}    % Divisiones y Ruffini

\usepackage{pgfplots}   %Representar funciones
\pgfplotsset{compat=1.18}  % Versión 1.18

%Definición de teoremas, etc.
\usepackage{amsthm}
%\swapnumbers   % Intercambia la posición del texto y de la numeración

\theoremstyle{plain}
\newtheorem{teo}{Teorema}[chapter]     % Se resetea en cada chapter
\newtheorem{coro}{Corolario}[teo]           % Se resetea en cada teorema
\newtheorem{prop}[teo]{Proposición}         % Usa el mismo contador que teorema
\newtheorem{lema}[teo]{Lema}                % Usa el mismo contador que teorema

\theoremstyle{remark}
\newtheorem*{observacion}{Observación}

\theoremstyle{definition}
\newtheorem{definicion}{Definición} [chapter]     % Se resetea en cada chapter
\newtheorem*{notacion}{Notacion}
\newtheorem*{ejemplo}{Ejemplo}
\newtheorem*{ejercicio*}{Ejercicio}             % No numerado
\newtheorem{ejercicio}{Ejercicio} [section]     % Se resetea en cada section


% Modificar el formato de la numeración del teorema "ejercicio"
\renewcommand{\theejercicio}{%
  \ifnum\value{section}=0 % Si no se ha iniciado ninguna sección
    \arabic{ejercicio}% Solo mostrar el número de ejercicio
  \else
    \thesection.\arabic{ejercicio}% Mostrar número de sección y número de ejercicio
  \fi
}


% \renewcommand\qedsymbol{$\blacksquare$}         % Cambiar símbolo QED
%------------------------------------------------------------------------

% Paquetes para encabezados
\usepackage{fancyhdr}
\pagestyle{fancy}
\fancyhf{}

\newcommand{\helv}{ % Modificación tamaño de letra
\fontfamily{}\fontsize{12}{12}\selectfont}
\setlength{\headheight}{15pt} % Amplía el tamaño del índice


%\usepackage{lastpage}   % Referenciar última pag   \pageref{LastPage}
\fancyfoot[C]{\thepage}

%------------------------------------------------------------------------

\begin{comment} % Descomentar en book. Comentar en article

% Conseguir que no ponga "Capítulo 1". Sino solo "1."
\makeatletter
\renewcommand{\@makechapterhead}[1]{%
  \vspace*{50\p@}%
  {\parindent \z@ \raggedright \normalfont
    \ifnum \c@secnumdepth >\m@ne
      \huge\bfseries \thechapter.\hspace{1em}\ignorespaces
    \fi
    \interlinepenalty\@M
    \Huge \bfseries #1\par\nobreak
    \vskip 40\p@
  }}
\makeatother
\end{comment}

%------------------------------------------------------------------------
% Personalizar la portada
\usepackage{titling}


%------------------------------------------------------------------------

% Comandos definidos
\newcommand{\bb}[1]{\mathbb{#1}}
\newcommand{\cc}[1]{\mathcal{#1}}
\newcommand{\Hop}{\stackrel{L'H\hat{o}pital}{=}}

\newcommand{\red}[1]{{\color{red}#1}} % Para integrales, destacar los cambios.

% Método de integración
\newcommand{\MetInt}[2]{
    \left[\begin{array}{c}
        #1 \\ #2
    \end{array}\right]
}

%------------------------------------------------------------------------


\begin{document}

    % 1. Foto de fondo
    % 2. Título
    % 3. Encabezado Izquierdo
    % 4. Color de fondo
    % 5. Coord x del titulo
    % 6. Coord y del titulo
    % 7. Fecha

    
    % 1. Foto de fondo
% 2. Título
% 3. Encabezado Izquierdo
% 4. Color de fondo
% 5. Coord x del titulo
% 6. Coord y del titulo
% 7. Fecha

\newcommand{\portada}[7]{

    \portadaBase{#1}{#2}{#3}{#4}{#5}{#6}{#7}
    \portadaBook{#1}{#2}{#3}{#4}{#5}{#6}{#7}
}

\newcommand{\portadaExamen}[7]{

    \portadaBase{#1}{#2}{#3}{#4}{#5}{#6}{#7}
    \portadaArticle{#1}{#2}{#3}{#4}{#5}{#6}{#7}
}




\newcommand{\portadaBase}[7]{

    % Tiene la portada principal y la licencia Creative Commons
    
    % 1. Foto de fondo
    % 2. Título
    % 3. Encabezado Izquierdo
    % 4. Color de fondo
    % 5. Coord x del titulo
    % 6. Coord y del titulo
    % 7. Fecha
    
    
    \thispagestyle{empty}               % Sin encabezado ni pie de página
    \newgeometry{margin=0cm}        % Márgenes nulos para la primera página
    
    
    % Encabezado
    \fancyhead[L]{\helv #3}
    \fancyhead[R]{\helv \nouppercase{\leftmark}}
    
    
    \pagecolor{#4}        % Color de fondo para la portada
    
    \begin{figure}[p]
        \centering
        \transparent{0.3}           % Opacidad del 30% para la imagen
        
        \includegraphics[width=\paperwidth, keepaspectratio]{assets/#1}
    
        \begin{tikzpicture}[remember picture, overlay]
            \node[anchor=north west, text=white, opacity=1, font=\fontsize{60}{90}\selectfont\bfseries\sffamily, align=left] at (#5, #6) {#2};
            
            \node[anchor=south east, text=white, opacity=1, font=\fontsize{12}{18}\selectfont\sffamily, align=right] at (9.7, 3) {\textbf{\href{https://losdeldgiim.github.io/}{Los Del DGIIM}}};
            
            \node[anchor=south east, text=white, opacity=1, font=\fontsize{12}{15}\selectfont\sffamily, align=right] at (9.7, 1.8) {Doble Grado en Ingeniería Informática y Matemáticas\\Universidad de Granada};
        \end{tikzpicture}
    \end{figure}
    
    
    \restoregeometry        % Restaurar márgenes normales para las páginas subsiguientes
    \pagecolor{white}       % Restaurar el color de página
    
    
    \newpage
    \thispagestyle{empty}               % Sin encabezado ni pie de página
    \begin{tikzpicture}[remember picture, overlay]
        \node[anchor=south west, inner sep=3cm] at (current page.south west) {
            \begin{minipage}{0.5\paperwidth}
                \href{https://creativecommons.org/licenses/by-nc-nd/4.0/}{
                    \includegraphics[height=2cm]{assets/Licencia.png}
                }\vspace{1cm}\\
                Esta obra está bajo una
                \href{https://creativecommons.org/licenses/by-nc-nd/4.0/}{
                    Licencia Creative Commons Atribución-NoComercial-SinDerivadas 4.0 Internacional (CC BY-NC-ND 4.0).
                }\\
    
                Eres libre de compartir y redistribuir el contenido de esta obra en cualquier medio o formato, siempre y cuando des el crédito adecuado a los autores originales y no persigas fines comerciales. 
            \end{minipage}
        };
    \end{tikzpicture}
    
    
    
    % 1. Foto de fondo
    % 2. Título
    % 3. Encabezado Izquierdo
    % 4. Color de fondo
    % 5. Coord x del titulo
    % 6. Coord y del titulo
    % 7. Fecha


}


\newcommand{\portadaBook}[7]{

    % 1. Foto de fondo
    % 2. Título
    % 3. Encabezado Izquierdo
    % 4. Color de fondo
    % 5. Coord x del titulo
    % 6. Coord y del titulo
    % 7. Fecha

    % Personaliza el formato del título
    \pretitle{\begin{center}\bfseries\fontsize{42}{56}\selectfont}
    \posttitle{\par\end{center}\vspace{2em}}
    
    % Personaliza el formato del autor
    \preauthor{\begin{center}\Large}
    \postauthor{\par\end{center}\vfill}
    
    % Personaliza el formato de la fecha
    \predate{\begin{center}\huge}
    \postdate{\par\end{center}\vspace{2em}}
    
    \title{#2}
    \author{\href{https://losdeldgiim.github.io/}{Los Del DGIIM}}
    \date{Granada, #7}
    \maketitle
    
    \tableofcontents
}




\newcommand{\portadaArticle}[7]{

    % 1. Foto de fondo
    % 2. Título
    % 3. Encabezado Izquierdo
    % 4. Color de fondo
    % 5. Coord x del titulo
    % 6. Coord y del titulo
    % 7. Fecha

    % Personaliza el formato del título
    \pretitle{\begin{center}\bfseries\fontsize{42}{56}\selectfont}
    \posttitle{\par\end{center}\vspace{2em}}
    
    % Personaliza el formato del autor
    \preauthor{\begin{center}\Large}
    \postauthor{\par\end{center}\vspace{3em}}
    
    % Personaliza el formato de la fecha
    \predate{\begin{center}\huge}
    \postdate{\par\end{center}\vspace{5em}}
    
    \title{#2}
    \author{\href{https://losdeldgiim.github.io/}{Los Del DGIIM}}
    \date{Granada, #7}
    \thispagestyle{empty}               % Sin encabezado ni pie de página
    \maketitle
    \vfill
}
    \portadaExamen{ffccA4.jpg}{Cálculo II\\Examen X}{Cálculo II. Examen X}{MidnightBlue}{-8}{28}{2023}

    \begin{description}
        \item[Asignatura] Cálculo II.
        \item[Curso Académico] 2016-17.
        \item[Grado] Doble Grado en Ingeniería Informática y Matemáticas.
        \item[Grupo] Único.
        %\item[Profesor] María Victoria Velasco Collado.
        \item[Descripción] Parcial 2.
        %\item[Fecha] 24 de abril de 2019.
        %\item[Duración] 60 minutos.
    
    \end{description}
    \newpage
    
    \begin{ejercicio} [\textbf{3 puntos}]
Enuncia con precisión el Teorema Fundamental del Cálculo. Demuéstralo.
    
\end{ejercicio}

\begin{ejercicio} [\textbf{3 puntos}]
    Decir si son verdaderas o falsas las siguientes cuestiones, justificando la respuesta:

    \begin{enumerate}
        \item Dada $f:\bb{R}\to \bb{R}$ cóncava hacia arriba, entonces $f\circ f$ es también cóncava hacia arriba.

        En el caso de que $f$ fuese creciente, tendríamos que es cierto (Ejercicio 13 de la Relación 3). Por tanto, para encontrar un contraejemplo, necesitamos buscar una función decreciente y cóncava hacia arriba.

        Sea $f(x)=e^{-x}$. Tenemos que es estrictamente decreciente. Veamos ahora que es cóncava hacia arriba.
        \begin{equation*}
            f'(x)=-e^{-x} \qquad f''(x)=e^{-x}>0\qquad \forall x\in \bb{R}
        \end{equation*}
        Por tanto, como $f''(x)>0\; \forall x\in \bb{R}$, tenemos que $f$ es cóncava hacia arriba.

        Veamos ahora si $f\circ f$ es cóncava hacia arriba. Como tenemos que $(f\circ f)\in C^\infty(\bb{R})$, usamos el criterio de la segunda derivada.
        \begin{equation*}
            (f\circ f)'(x)=f'(f(x))f'(x) = -e^{-e^{-x}}(-e^{-x}) = e^{-e^{-x}}e^{-x} = e^{-e^{-x}-x}
        \end{equation*}
        \begin{equation*}
            (f\circ f)''(x)=e^{-e^{-x}-x} (e^{-x}-1)
        \end{equation*}

        Para $x>0$, tenemos que $(f\circ f)''(x)<0$, por lo que $f\circ f$ es cóncava hacia abajo en $\bb{R}^+$. Por tanto, el enunciado es \textbf{falso}.
        
    
        \item Sean $f:A\to \bb{R},g:B\to \bb{R}$ con $f(A)\subset B$ dos funciones uniformemente continuas. Entonces $g\circ f$ es uniformemente continua.

        Tenemos que $f$ es uniformemente continua; es decir:
        \begin{equation*}
            \forall \hat{\varepsilon} > 0, \exists \delta>0 \mid \text{Si } x,y\in A, \text{ con } |x-y|<\delta \Longrightarrow |f(x)-f(y)|<\hat{\varepsilon}
        \end{equation*}

        Tenemos que $g$ es uniformemente continua; es decir:
        \begin{equation*}
            \forall {\varepsilon} > 0, \exists \hat{\delta}>0 \mid \text{Si } x,y\in B, \text{ con } |x-y|<\hat{\delta} \Longrightarrow |g(x)-g(y)|<{\varepsilon}
        \end{equation*}

        En particular, tomando $\hat{\varepsilon} = \hat{\delta}$, y usando que $f(A)\subset B$, la continuidad uniforme de $g$ queda:
        \begin{equation*}
            \forall {\varepsilon} > 0, \exists \hat{\varepsilon}>0 \mid \text{Si } f(x),f(y)\in B, \text{ con } |f(x)-f(y)|<\hat{\varepsilon} \Longrightarrow |g(f(x))-g(f(y))|<{\varepsilon}
        \end{equation*}

        Uniendo lo que tenemos, queda:
        \begin{equation*}
            \forall \varepsilon>0, \exists \delta>0\mid  \text{Si } x,y\in A, \text{ con } |x-y|<\delta \Longrightarrow |(g\circ f)(x)-(g\circ f)(y)|<{\varepsilon}
        \end{equation*}

        Es decir, se ha demostrado que $g\circ f$ es uniformemente continua, por lo que es \textbf{cierto}.

        \item Toda función $f:\bb{R}^+_0\to \bb{R}$ uniformemente continua está acotada.

        Tomando $f(x)=x$, tenemos que es lipsitchziana por ser $f'(x)=1$ acotada. Al ser $f$ lipsitchziana, tenemos que $f$ es uniformemente continua.

        No obstante, $f$ no está acotada, por lo que el enunciado es \textbf{falso}.

        \item Si $f:\bb{R}^+\to \bb{R}$ es uniformemente continua y tiene límite en $+\infty$, es una función acotada.

        Por reducción al absurdo, supongamos que $f$ no está acotada. Entonces, $$\exists \{x_n\}\;(x_n\in \bb{R}^+\;\forall n\in \bb{N}) \mid \{|f(x_n)|\}\to +\infty$$

        Como $f$ tiene límite en $+\infty$, se tiene que: $$\forall \{x_n'\}\to +\infty\;(x_n'\in \bb{R}^+\;\forall n\in \bb{N}) \Longrightarrow \{f(x_n')\}\to L\in \bb{R}$$

        Por tanto, tenemos que $\{x_n\}$ no diverge, por lo que está acotada.
        \begin{equation*}
            \exists M>0 \mid |x_n|\leq M \qquad \forall n\in \bb{N}
        \end{equation*}

        Consideramos ahora $f_{]0,M[}$. Como $f$ es uniformemente continua, transforma conjuntos acotados en conjuntos acotados. Por tanto, $f(]0,M[)$ está acotado.

        Por tanto, tenemos que $f(x_n)\in f(]0,M[) \quad \forall n\in \bb{N}$. Por tanto, $f(x_n)$ acotado. Pero como $\{|f(x_n)|\}$ diverge, llegamos a una contradicción.

        Por tanto, $f$ está acotada.

        \begin{observacion}
            Tenemos que $\forall \{x_n\}\to +\infty$, tenemos que la imagen de la cola está acotada por $L+\varepsilon$, y que la imagen de los primeros términos, por ser un conjunto acotado y ser $f$ uniformemente continua, también es uniformemente continua.
        \end{observacion}
    \end{enumerate}
\end{ejercicio}

\begin{ejercicio} [\textbf{4 puntos}]
    Definimos $G:]-1,1[ \to \bb{R}$ como $G(x)=\int_0^{x^2} \frac{t^2dt}{\sqrt{1-t^2}}$.

    \begin{enumerate}
        \item Calcular la imagen de $G$.

        Como el integrando es una función continua en el intervalo de definición, tenemos que es Riemman Integrable. Por el Teorema Fundamental del Cálculo, tenemos que es derivable en $]-1,1[$ con derivada:
        \begin{equation*}
            G'(x)=\frac{x^4}{\sqrt{1-x^4}}\cdot 2x
            = \frac{2x^5}{\sqrt{1-x^4}} = 0 \Longleftrightarrow x=0
        \end{equation*}

        Por tanto, el único punto crítico es $x=0$. Tenemos $G(0)=0$. Además, como $\forall x^2\in ]0,1[$ tenemos que la imagen del integrando es positiva, tenemos que $G(x)\geq 0\;\forall x\in ]-1,1[$. Por tanto, se trata de un mínimo relativo. Al ser $G$ continua y solo tener un punto que anula a la primera derivada, es también absoluto. 

        Calculamos las siguientes imágenes:
        \begin{multline*}
            \lim_{x\to 1}G(x)=\int_0^1 \frac{t^2}{\sqrt{1-t^2}}dt
            = \lim_{b\to 1}\int_0^b \frac{t^2}{\sqrt{1-t^2}}dt
            =\MetInt{\sen u =t \quad u\in \left[0,\frac{\pi}{2}\right]}{\cos u\; du = dt}
            =\\=\lim_{b\to \frac{\pi}{2}}\int_0^b \frac{\sen^2 u}{\sqrt{1-\sen ^2 u}}\cos u \;du
            =\int_0^{\frac{\pi}{2}} \frac{\sen^2 u}{\sqrt{\cos ^2 u}}\cos u \;du
        \end{multline*}

        Por el intervalo empleado para el cambio de variable, tengo que $\cos u >0$. Por tanto,
        \begin{multline*}
            \lim_{x\to 1}G(x)=\lim_{b\to \frac{\pi}{2}}\int_0^b \frac{\sen^2 u}{\cos u}\cos u \;du
            = \lim_{b\to \frac{\pi}{2}}\int_0^b \sen^2 u \;du
            = \lim_{b\to \frac{\pi}{2}}\int_0^b \frac{1-\cos(2t)}{2} \;du
            =\\=\lim_{b\to \frac{\pi}{2}} \left[\frac{t-\sen(2t)}{2}\right]_0^b
            = \frac{\pi}{4}
        \end{multline*}

        Como además $G(-x)=G(x)$ por ser $G$ una función par, tenemos que:
        \begin{equation*}
            Im(G)=\left[0,\frac{\pi}{4} \right[
        \end{equation*}
        
        \item Decidir si tiene solución la ecuación $G(x)=\frac{1}{3}$. ¿Y la ecuación $G(x)=2$?

        $G(x)$ es continua y definida en un intervalo. Tenemos que $Im(G)=\left[0,\frac{\pi}{4} \right[$. Por tanto,
        \begin{equation*}
            \exists x\in ]-1,1[\mid G(x)=\frac{1}{3}\Longleftrightarrow \frac{1}{3}\in Im(G) \Longleftrightarrow \frac{1}{3} < \frac{\pi}{4} \Longleftrightarrow 4<3\pi  \qquad \text{Cierto}
        \end{equation*}
        \begin{equation*}
            \exists x\in ]-1,1[\mid G(x)=2\Longleftrightarrow 2\in Im(G) \Longleftrightarrow 2 < \frac{\pi}{4} \Longleftrightarrow 8<\pi \qquad \text{Falso}
        \end{equation*}

        Por tanto, tenemos que para la primera ecuación sí hay solución, pero para la segunda no.
    \end{enumerate}
\end{ejercicio}
    



\end{document}